\documentclass[11pt, twoside]{article}
\usepackage{amssymb}
\usepackage{microtype} 
\usepackage{amsfonts}
\usepackage{amsmath}
\usepackage{anysize}
\usepackage{amsthm}
\usepackage{color}
\usepackage{mathrsfs}
\usepackage{txfonts}
\usepackage{latexsym}
\usepackage{indentfirst}
\usepackage{bm}
\usepackage{url}
\urlstyle{same}

%\usepackage{refcheck}
%\usepackage{showkeys}

\usepackage[colorlinks=true,
linkcolor=red,
citecolor=blue,
urlcolor=magenta]{hyperref}
%\usepackage{refcheck}
\allowdisplaybreaks


\numberwithin{equation}{section}

\pagestyle{myheadings}
\markboth{\footnotesize\rm\sc Ziming Wang and Ziyi He}
{\footnotesize\rm\sc Weighted sixth moments on residue grids}



\textwidth=15cm
\textheight=24.129cm
\oddsidemargin 0.46cm
\evensidemargin 0.46cm


\newtheorem{theorem}{Theorem}[section]
\newtheorem{lemma}[theorem]{Lemma}
\newtheorem{corollary}[theorem]{Corollary}
\newtheorem{proposition}[theorem]{Proposition}
\newtheorem{claim}[theorem]{Claim}
\newtheorem{example}[theorem]{Example}
\newtheorem{problem}[theorem]{Open Problem}
\theoremstyle{definition}
\newtheorem{remark}[theorem]{Remark}
\newtheorem{definition}[theorem]{Definition}
\newtheorem{assumption}[theorem]{Assumption}
\renewcommand{\appendix}{\par
	\setcounter{section}{0}
	\setcounter{subsection}{0}
	\setcounter{subsubsection}{0}
	\gdef\thesection{\@Alph\c@section}
	\gdef\thesubsection{\@Alph\c@section.\@arabic\c@subsection}
	\gdef\theHsection{\@Alph\c@section.}
	\gdef\theHsubsection{\@Alph\c@section.\@arabic\c@subsection}
	\csname appendixmore\endcsname
}

\renewcommand{\theenumi}{\roman{enumi}}
\renewcommand{\labelenumi}{\rm (\theenumi)}
\newcommand{\T}{\mathbb T}
\newcommand{\Z}{\mathbb Z}
\newcommand{\C}{\mathbb C}
\newcommand{\E}{\mathbb E}
\newcommand{\cE}{\mathcal E}
\newcommand{\cG}{\mathcal G}
\newcommand{\ind}{\mathbf 1}
\newcommand{\dd}{\,\mathrm d}
\newcommand{\Lcal}{\mathfrak L}
\newcommand{\Ksix}{\mathfrak K_6}
\DeclareMathOperator{\supp}{supp}
\begin{document}
\title{\bf\Large Weighted Sixth Moments for the Parabola on Rectangular Residue Grids
\footnotetext{\hspace{-0.35cm} 2020 {\it Mathematics Subject Classification}. Primary 42B10; Secondary 11L07, 11B30, 42B37.\endgraf
{\it Key words and phrases.}\ discrete restriction, sixth moment, composite modulus, parabola,
modular line.\endgraf
This project is supported by the National Natural Science Foundation of China (Grant No.\ 12201061)
and the Fundamental Research Funds for the Central Universities (Grant No.\ 2023ZCJH02).}}
\author{Ziming Wang\footnote{Corresponding author,
E-mail: \texttt{zimingwang@bupt.edu.cn}/{\color{red}\today}/Final version.}
\ and Ziyi He}
\date{}
\maketitle	
\vspace{-0.8cm}
	
\begin{center}
\begin{minipage}{13cm}
{\small {\bf Abstract}\quad   
Let $N\geq 2$ and $N\leq M\leq N^2$.  On the probability space
$(\Z/N\Z)\times(\Z/M\Z)$, consider the parabolic extension operator
$$
\cE_{N,M}a(A,B)=\sum_{n=0}^{N-1}a_n e\!\left(\frac{nA}{N}+\frac{n^2B}{M}\right).
$$
We prove the weighted sixth-moment estimate
$$
\E_{A\bmod N}\E_{B\bmod M}|\cE_{N,M}a(A,B)|^6
\leq C\left\{N^2\log(2N)+\frac{N^4}{M}\Lcal_{<N}(M)\right\}
\|a\|_{\ell^6}^6,
$$
where
$$
\Lcal_{<N}(M):=\sum_{\substack{d\mid M\\d<N}}\frac{\varphi(d)}{d^2}.
$$
For $q\ge 6$, the setimate is a consequence of H\"older's inequality; for $2\le q<6$, interpolation between $\ell^2$ and $\ell^6$
endpoints gives the results.}
\end{minipage}
\end{center}


\section{Introduction}
Let $e(t):=\exp(2\pi i t)$ and $I_N=\{0,\ldots,N-1\}$. For $a:I_N\to\C$, set
$$
\cE_{N,M}a(A,B)
:=
\sum_{n=0}^{N-1}a_n e\!\left(\frac{nA}{N}+\frac{n^2B}{M}\right), \quad \E H=\frac1{NM}\sum_{A=0}^{N-1}\sum_{B=0}^{M-1}H(A,B),
$$
and equip $(\Z/N\Z)\times(\Z/M\Z)$ with normalized counting measure.
We study the best weighted constant
$$
\Ksix(N,M)
:=
\sup_{a\neq0}\frac{
\E |\cE_{N,M}a|^6}
{\|a\|_{\ell^6(I_N)}^6}.
$$

Hickman and Wright \cite{HickmanWright2018} developed restriction and Kakeya theory in $\Z/N\Z$.  They gives the diagonal estimate
$$
\E_{A,B\bmod N}|\cE_{N,N}a|^6
\lesssim N\|a\|_{\ell^2}^6
\lesssim N^3\|a\|_{\ell^6}^6,
\qquad N=p^\alpha,
$$
with $p$ odd. Johnsrude \cite{Johnsrude2024} obtained the estimate $N^2+N^4/M$ for a family
of common odd prime powers by small cap decoupling. His argument gives
$$
	\E
	\left|\cE_{N,M}a\right|^6
	\lesssim_p
	(\log N)^{23+6\sigma}
	\left(N^2+\frac{N^4}{M}\right)
	\sum_{n=0}^{N-1}|a_n|^6,
$$
when $N=p^K$ and $M=N^{1+\sigma}$ are powers of the same odd prime. Johnsrude \cite{Johnsrude2025Sparse} also compares
$p$-adic local restriction constants with real sparse-domain constants and applies the comparison
to parabolic sparse strips.
Kingsbury-Neuschotz \cite{KingsburyNeuschotz2025} treats the diagonal squarefree restriction problem. Moreover, Cushman et
al. \cite{CushmanDemeterWu2026} prove
$$
J_3(\gamma(X))\ll_\varepsilon |X|^{3+\varepsilon}
$$
for every finite subset of a strictly convex graph
.  Applying this hereditary
estimate to dyadic level sets on the parabola gives the weighted bound
$$
\E|\cE_{N,M}a|^6\lesssim_\varepsilon
N^\varepsilon\left(1+\frac{N^2}{M}\right)\|a\|_2^6
\lesssim_\varepsilon N^\varepsilon\left(N^2+\frac{N^4}{M}\right)\|a\|_6^6.
$$
Besides, Hu’s \cite{Hu2026} turning-complexity extension gives the same unweighted input for the parabola, where the
turning complexity equals one.

Thus, the results above inspire us to consider the question: Is it possible to establish a sharp sixth-moment estimate for the
parabolic extension operator?

\begin{theorem}\label{thm:main-q6}
Let $N,M\in\Z$ satisfy $N\ge 2$ and $N\leq M\leq N^2$. Then, for every $a:I_N\to\C$,
$$
\E|\cE_{N,M}a|^6
\leq
C\left\{
N^2\log(2N)+\frac{N^4}{M}\Lcal_{<N}(M)
\right\}\|a\|_{\ell^6}^6,
$$
where $C$ is absolute and
$$
\Lcal_{<N}(M)
:=
\sum_{\substack{d\mid M\\d<N}}\frac{\varphi(d)}{d^2}.
$$
\end{theorem}

\begin{corollary}\label{cor:uniform-q6}
Under the assumptions of Theorem~\ref{thm:main-q6},
\begin{equation}\label{eq:loglog-uniform}
\E|\cE_{N,M}a|^6
\leq
C\left\{
N^2\log(2N)
+\frac{N^4}{M}\bigl(1+\log\log(3M)\bigr)
\right\}\|a\|_{\ell^6}^6.
\end{equation}
If $M=p^\alpha$ is a prime power, then
\begin{equation}\label{eq:prime-power-branch}
\E|\cE_{N,M}a|^6
\leq
C\left\{N^2\log(2N)+\frac{N^4}{M}\right\}\|a\|_{\ell^6}^6.
\end{equation}
In particular, a rough estimate
\begin{equation}\label{eq:common-log}
\E|\cE_{N,M}a|^6
\leq
C\log(2N)\left(N^2+\frac{N^4}{M}\right)\|a\|_{\ell^6}^6
\end{equation}
holds uniformly.
\end{corollary}
Applying the standard Riesz-Thorin interpolation theorem to the estimates above yields the corresponding mixed-norm bounds.
To streamline the statement of the following theorem, we introduce the shorthand notations
$$
L_N:=\log(2N),
\qquad
X:=\frac{N^2}{M},
\qquad
D_{N,M}:=L_N+X\Lcal_{<N}(M).
$$
With these definitions at hand, we have the following mixed-norm refinement.
\begin{theorem}\label{thm:mixed-q}
Suppose that $2\leq q\le\infty$. Let $\theta=\frac{3}{2}-\frac{3}{q}$, then 
\begin{equation}\label{eq:mixed-refined}
\E|\cE_{N,M}a|^6
\leq
C_{q}
N^{3-6/q}
L_N^{(1-\theta)}
(1+X)^{1-\theta}
D_{N,M}^{\theta}
\|a\|_{\ell^q}^6.
\end{equation}
A rough estimate gives
\begin{equation}\label{eq:mixed-clean}
\E|\cE_{N,M}a|^6
\leq
C_{q}
L_N
\left(N^{3-6/q}+\frac{N^{5-6/q}}{M}\right)
\|a\|_{\ell^q}^6.
\end{equation}
Moreover, if $6\leq q\leq\infty$, then
\begin{equation}\label{eq:mixed-qge6}
\E|\cE_{N,M}a|^6\leq C N^{3-6/q}D_{N,M}\|a\|_{\ell^q}^6.
\end{equation}
\end{theorem}

This paper is organized as follows. In Section~\ref{sec:preliminaries}, we introduce some basic 
definitions and useful lemmas.
In Section~\ref{sec:main-proof}, we prove Theorem~\ref{thm:main-q6}.
In Section~\ref{sec:arithmetic-factor}, we study the arithmetic factor and prove
Corollary~\ref{cor:uniform-q6}.
In Section~\ref{sec:lower-bounds}, we give the lower bounds and classify the
sharp regimes.
In Section~\ref{sec:interpolation}, we prove Theorem~\ref{thm:mixed-q}.
In Section~\ref{sec:continuous}, we deduce the continuous cell estimate.
And Section~\ref{sec:conclusion} states the remaining question.

We end this introduction by making some conventions on the notation. Let $\Z:=\{1,2,\cdots\}$. We denote by $C$ a positive constant
which is independent of the main parameters, 
but may vary from line to line. We use $C_{(\alpha,\cdots)}$ to denote a positive constant depending on the indicated parameters
$\alpha,\ldots$.  The symbol $f\lesssim g$ means $f\le C g$ and, if $f\lesssim g \lesssim f$, then we write $f\sim g$.
\section{Preliminaries}\label{sec:preliminaries}

For convenience, we begin by fixing the notation used throughout the paper. We set 
$N,M\in\Z$, $N\geq2$, $N\leq M\leq N^2$. 
Let \(I_N \coloneqq \{0,\ldots,N-1\}\), and write 
\(e_Q(z) \coloneqq \exp(2\pi i z/Q)\) for \(Q\ge 1\). 
The group \(G_{N,M} \coloneqq (\mathbb{Z}/N\mathbb{Z}) \times (\mathbb{Z}/M\mathbb{Z})\) 
is equipped with the normalized counting measure, so that for any function \(H\) on \(G_{N,M}\),
$$
\mathbb{E} H = \frac{1}{NM}\sum_{A=0}^{N-1}\sum_{B=0}^{M-1} H(A,B).
$$
Let $a:I_N\rightarrow\mathbb{C}$. For \(1\le q\le \infty\), the \(\ell^q\)-norm of $a$ is defined as
$$
\|a\|_{\ell^q}:=\left(\sum_{n=0}^{N-1}|a(n)^q|\right)^{1/q},
$$
and for $q=\infty$,
$$
\|a\|_{\ell^\infty}:=\max_{0\le n\le N-1}|a(n)|.
$$
With these notations, the parabolic extension operator takes the form
$$
\mathcal{E}_{N,M}a(A,B) = \sum_{n\in I_N} a_n e_N(nA) e_M(n^2B).
$$
A standard expansion of the sixth moment, combined with finite character orthogonality, gives the exact combinatorial identity
\begin{equation}\label{eq:six-moment-expansion}
\mathbb{E}|\mathcal{E}_{N,M}a|^6 = 
\sum_{(n,m)\in \mathcal{G}_{N,M}} 
a_{n_1}a_{n_2}a_{n_3} \overline{a_{m_1}a_{m_2}a_{m_3}},
\end{equation}
where \(\mathcal{G}_{N,M}\) denotes the set of ordered sextuples \((n,m) \in I_N^3\times I_N^3\) satisfying
\begin{equation}\label{eq:G-def}
\cG_{N,M}:=\left\{(\bm n,\bm m)\in I_N^3\times I_N^3:
\begin{array}{l}
n_1+n_2+n_3\equiv m_1+m_2+m_3\pmod N,\\
n_1^2+n_2^2+n_3^2\equiv m_1^2+m_2^2+m_3^2\pmod M
\end{array}
\right\}.
\end{equation}
Thus, the two terms in Theorem \ref{thm:main-q6} will arise respectively from the linear congruence and the quadratic congruence.
In carrying out this counting, the following arithmetic divisor sum naturally appears:
$$
\mathfrak{L}_{<N}(M) \coloneqq \sum_{\substack{d\mid M\\ d<N}} \frac{\phi(d)}{d^2},
\qquad\text{and}\qquad
\mathfrak{L}(M) \coloneqq \sum_{d\mid M} \frac{\phi(d)}{d^2}.
$$

Now, we show some technical lemmas which are used throughout this paper. The following lemma comes 
from some simple estimates. We omit the details here.
\begin{lemma}[{\cite[Theorems 2.7 and 6.9]{MontgomeryVaughan2006}}]\label{estimate}
Let $p$ denote a prime number. For any $y \ge 2$, the following estimates hold:
$$
\prod_{p \le y} \left(1 - \frac{1}{p}\right)^{-1} \lesssim \log(2y),
$$
and
$$
\vartheta(y) := \sum_{p \le y} \log p \sim y \qquad (y \to \infty).
$$
\end{lemma}

\begin{lemma}[{\cite[Theorem~1.2]{DFGGL26}}]\label{Guo}
Let $S\subset \mathbb{Z}$ be a finite set with $\# S\ge 2$, one has 
\begin{equation}\label{eq:gly-input}
\left\|\sum_{n=1}^{N}b_ne(nx+n^2t)\right\|_{L^6(\T^2)}\leq C_6\bigl(\log\# S\bigr)^{1/6}
\left(\sum_{n=1}^{N}|b_n|^2\right)^{1/2}.
\end{equation}
\end{lemma}\label{J}
Next lemma is just in \cite{Rogovskaya1986} (see also in \cite{VaughanWooley1997}).
\begin{lemma}
Let $J_{3,2}(N)$ denote the number of $(\bm x,\bm y)\in\{1,\ldots,N\}^6$ satisfying
$$
x_1+x_2+x_3=y_1+y_2+y_3\quad\text{and}\quad
x_1^2+x_2^2+x_3^2=y_1^2+y_2^2+y_3^2.
$$

\begin{equation}\label{eq:J-asymptotic}
J_{3,2}(N)=\frac{18}{\pi^2}N^3\log N+O(N^3).
\end{equation}
\end{lemma}

\begin{lemma}\label{lem:integer-line}
Let $(u,v)\in\Z^2\setminus\{(0,0)\}$ and put
$$
g:=\gcd(|u|,|v|) \quad\text{and}\quad H:=\max(|u|,|v|).
$$
For every $c\in\Z$,
$$
\#\{(x,y)\in I_N^2:ux+vy=c\}\leq1+\frac{Ng}{H}.
$$
\end{lemma}

\begin{proof}
If the equation has no integer solution, the bound is immediate.
Otherwise, fix an integer solution $(x_0,y_0)$.  Every integer solution has
the form
$$
(x,y)=\left(x_0+\frac vg s,y_0-\frac ug s\right),
\qquad s\in\Z.
$$
At least one coordinate changes by $H/g$ between consecutive parameter
values.  That coordinate must lie in $I_N$, an interval containing $N$
consecutive integers.  Hence there are at most $1+Ng/H$ admissible values of
$s$.  This parametrization also covers $u=0$ or $v=0$.
\end{proof}

\begin{lemma}\label{lem:primitive-sum}
There exists a constant $C$ such that, for every $R\geq1$,
$$
\sum_{\substack{(u,v)\in\Z^2\\0<\max(|u|,|v|)\leq R}}
\frac{\gcd(|u|,|v|)}{\max(|u|,|v|)}\leq CR\log(2R).
$$
\end{lemma}

\begin{proof}
Write every nonzero direction uniquely as
$$
(u,v)=g(u_0,v_0),\qquad g\geq1,\qquad\gcd(|u_0|,|v_0|)=1.
$$
For $s=\max(|u_0|,|v_0|)$, the square boundary
$\max(|u_0|,|v_0|)=s$ contains at most $8s$ integer points.
The allowed multiples satisfy $1\leq g\leq R/s$, and each contributes $1/s$.
Thus
$$
\sum_{0<\max(|u|,|v|)\leq R}\frac{\gcd(|u|,|v|)}{\max(|u|,|v|)}
\leq\sum_{s=1}^{R}(8s)\left\lfloor\frac Rs\right\rfloor\frac1s
\leq 8R\sum_{s=1}^{R}\frac1s,
$$
which proves the claim.
\end{proof}

\begin{lemma}\label{lem:modular-line}
Let $(u,v)\in\Z^2\setminus\{(0,0)\}$ with $|u|,|v|<N$, and put
$$
g:=\gcd(|u|,|v|),\qquad H:=\max(|u|,|v|),\qquad\delta:=\gcd(2g,M).
$$
For every $c\in\Z$,
$$
\#\{(x,y)\in I_N^2:2ux+2vy\equiv c\pmod M\}\leq C\left(1+\frac{Ng}{H}+\frac{N^2\delta}{M}\right).
$$
\end{lemma}

\begin{proof}
Each solution of the congruence determines a unique $\ell\in\Z$ such that
$$
2ux+2vy=c+\ell M.
$$
For $(x,y)$ in $I_N^2$, the left-hand side lies in an interval of length at
most
$$
2(|u|+|v|)(N-1)\leq4HN.
$$
Hence the possible values of $\ell$ lie in an interval of length at most
$4HN/M$.  The corresponding line contains an integer point if and only if
$$
2g\mid c+\ell M.
$$
If $\ell_0$ gives one such line, then
$$
2g\mid M(\ell-\ell_0)
\quad\Longleftrightarrow\quad
\frac{2g}{\delta}\mid\ell-\ell_0.
$$
The number of lifts containing an integer point is therefore at most
$$
1+\frac{4HN}{M}\frac{\delta}{2g}
=1+\frac{2HN\delta}{Mg}.
$$
Using Lemma~\ref{lem:integer-line} gives at most $1+Ng/H$
box points on each lift.  Multiplying the two bounds gives
$$
\left(1+\frac{2HN\delta}{Mg}\right)
\left(1+\frac{Ng}{H}\right).
$$
The terms other than $1+Ng/H$ are
$$
\frac{2HN\delta}{Mg}+\frac{2N^2\delta}{M}.
$$
Since $H/g\leq H<N$, the first term is at most $2N^2\delta/M$.
This finishes the proof of Lemma \ref{lem:modular-line}. 
\end{proof}

\begin{lemma}\label{lem:gcd-direction}
Let
$
\mathcal D_N:=\{(u,v)\in\Z^2:|u|,|v|<N,(u,v)\neq(0,0)\}.
$
Then we have
$$
\sum_{(u,v)\in\mathcal D_N}
\gcd\bigl(2\gcd(|u|,|v|),M\bigr)
\leq CN^2\Lcal_{<N}(M).
$$
\end{lemma}

\begin{proof}
We set
$$
r(u,v):=\gcd(|u|,|v|,M).
$$
If $g=\gcd(|u|,|v|)$, then
$$
\gcd(2g,M)\leq2\gcd(g,M)=2r(u,v).
$$
Using $r=\sum_{d\mid r}\varphi(d)$ and interchanging finite sums gives
\begin{equation}\label{eq:gcd-expand}
\begin{aligned}
\sum_{(u,v)\in\mathcal D_N}r(u,v)
&=\sum_{d\mid M}\varphi(d)
\#\{(u,v)\in\mathcal D_N:d\mid u,\ d\mid v\}\\
&=\sum_{d\mid M}\varphi(d)
\left\{\left(2\left\lfloor\frac{N-1}{d}\right\rfloor+1\right)^2-1\right\}.
\end{aligned}
\end{equation}
The expression in braces vanishes for $d\geq N$.  If $d<N$ and
$q=\lfloor(N-1)/d\rfloor\geq1$, then
$$
(2q+1)^2-1=4q(q+1)\leq8q^2\leq\frac{8N^2}{d^2}.
$$
Substitution into \eqref{eq:gcd-expand}, followed by
$\gcd(2g,M)\leq2r(u,v)$, proves the lemma.
\end{proof}

\section{Proof of the sixth-moment estimate}\label{sec:main-proof}

We first bound the number of completions of a fixed coordinate in
\eqref{eq:G-def}.  This bound will then give the estimate for arbitrary
weights.

\begin{proposition}\label{prop:fixed-slot}
Fix any one of the six coordinates in $\cG_{N,M}$ and prescribe its value
$j\in I_N$.  The number of admissible completions is at most
\begin{equation}\label{eq:fixed-slot}
C\left\{N^2\log(2N)+\frac{N^4}{M}\Lcal_{<N}(M)\right\}.
\end{equation}
\end{proposition}

\begin{proof}
The admissible set is invariant under permutations within each triple and
under interchange of $\bm n$ and $\bm m$.  It therefore suffices to fix
$n_1=j$.  For an admissible sextuple, set
$$
d_i:=n_i-m_i.
$$
Since $|d_i|<N$, the linear congruence uniquely determines
$$
h\in\{-2,-1,0,1,2\},
\qquad
d_1+d_2+d_3=hN.
$$
Indeed, the sum of the three differences has absolute value less than $3N$
and is divisible by $N$.

Fix $h$ and write
$$
u:=d_2,
\qquad
v:=d_3,
\qquad
d_1:=hN-u-v.
$$
Substituting $m_i=n_i-d_i$ into the quadratic congruence gives
$$
2un_2+2vn_3
\equiv d_1^2+u^2+v^2-2jd_1\pmod M.
$$
Once $j,h,u,v,n_2,n_3$ are fixed, the negative coordinates are determined by
$$
m_1=j-d_1,
\qquad
m_2=n_2-u,
\qquad
m_3=n_3-v.
$$
Dropping the conditions that these integers belong to $I_N$ gives an upper
bound.

Suppose first that $(u,v)\neq(0,0)$, and set
$$
g=\gcd(|u|,|v|),
\qquad
H=\max(|u|,|v|),
\qquad
\delta(u,v)=\gcd(2g,M).
$$
By Lemma~\ref{lem:modular-line}, the number of completions in this direction
is at most
$$
C\left(1+\frac{Ng}{H}+\frac{N^2\delta(u,v)}{M}\right).
$$
Summing over nonzero directions, the constant term contributes $O(N^2)$.
Lemma~\ref{lem:primitive-sum} gives
$$
N\sum_{(u,v)\in\mathcal D_N}\frac{g}{H}
\leq CN^2\log(2N),
$$
and Lemma~\ref{lem:gcd-direction} gives
$$
\frac{N^2}{M}
\sum_{(u,v)\in\mathcal D_N}\delta(u,v)
\leq C\frac{N^4}{M}\Lcal_{<N}(M).
$$

If $(u,v)=(0,0)$, then $d_1=hN$.  The condition $|d_1|<N$ forces
$h=d_1=0$, and hence $m_i=n_i$.  With $n_1=j$ fixed, there are exactly
$N^2$ completions.  This contribution is absorbed by the first term in
\eqref{eq:fixed-slot}.
\end{proof}

\begin{proof}[Proof of Theorem~\ref{thm:main-q6}]
Apply the triangle inequality to \eqref{eq:six-moment-expansion}.
For the six absolute values in each summand, use the AM--GM inequality
$$
\prod_{\nu=1}^{6}z_\nu
\leq\frac16\sum_{\nu=1}^{6}z_\nu^6,
\qquad z_\nu\geq0.
$$
Let $\deg_s(j)$ be the number of completions when the $s$th coordinate is
fixed at $j$, and let
$$
D^*_{N,M}:=
C\left\{N^2\log(2N)+\frac{N^4}{M}\Lcal_{<N}(M)\right\}.
$$
Proposition~\ref{prop:fixed-slot} gives
$\deg_s(j)\leq D^*_{N,M}$ for every $s\in\{1,\ldots,6\}$ and $j\in I_N$.
Consequently,
$$
\begin{aligned}
\E|\cE_{N,M}a|^6
&\leq\frac16\sum_{s=1}^{6}\sum_{j\in I_N}|a_j|^6\deg_s(j)\\
&\leq D^*_{N,M}\sum_{j\in I_N}|a_j|^6.
\end{aligned}
$$
This finishes the proof of Lemma \ref{thm:main-q6}.  The summands in
\eqref{eq:six-moment-expansion} may be complex.  Taking absolute values
before applying AM--GM justifies the argument for arbitrary complex weights.
\end{proof}

\section{The arithmetic factor and Corollary~\ref{cor:uniform-q6}}
\label{sec:arithmetic-factor}

\begin{proposition}\label{prop:euler-factor}
For every $N\geq2$ and every positive integer $M$,
\begin{equation}\label{eq:euler-product}
\Lcal(M)
=\prod_{p^\alpha\parallel M}
\left(1+\frac1p-\frac1{p^{\alpha+1}}\right).
\end{equation}
There is an absolute constant $C$ such that
\begin{equation}\label{eq:L-loglog}
1\leq\Lcal_{<N}(M)\leq\Lcal(M)
\leq C\bigl(1+\log\log(3M)\bigr).
\end{equation}
If $M=p^\alpha$, then
\begin{equation}\label{eq:L-prime-power}
\Lcal(M)=1+\frac1p-\frac1{p^{\alpha+1}}\leq\frac32.
\end{equation}
\end{proposition}

\begin{proof}
The function $d\mapsto\varphi(d)/d^2$ is multiplicative.  For
$p^\alpha\parallel M$, its local sum is
$$
\begin{aligned}
\sum_{j=0}^{\alpha}\frac{\varphi(p^j)}{p^{2j}}
&=1+\sum_{j=1}^{\alpha}\frac{p^j-p^{j-1}}{p^{2j}}\\
&=1+\left(1-\frac1p\right)\sum_{j=1}^{\alpha}\frac1{p^j}\\
&=1+\frac1p-\frac1{p^{\alpha+1}}.
\end{aligned}
$$
Multiplying the local sums proves \eqref{eq:euler-product}.
The prime-power bound \eqref{eq:L-prime-power} follows.

To proof \eqref{eq:L-loglog}, we first note that
$$
\Lcal(M)\leq\prod_{p\mid M}\left(1+\frac1p\right).
$$
Set $Y=\log(3M)$.  By using Lemma \ref{estimate}, we have that
$$
\prod_{\substack{p\mid M\\p\leq Y}}
\left(1+\frac1p\right)
\leq
\prod_{p\leq Y}\left(1-\frac1p\right)^{-1}
\leq C\log(2Y).
$$
If $M$ has $r$ distinct prime divisors greater than $Y$, then $Y^r\leq M$,
so
$$
r\leq\frac{\log M}{\log Y}.
$$
It follows that
$$
\prod_{\substack{p\mid M\\p>Y}}
\left(1+\frac1p\right)
\leq
\exp\!\left(\sum_{\substack{p\mid M\\p>Y}}\frac1p\right)
\leq
\exp\!\left(\frac{\log M}{Y\log Y}\right)
\leq C.
$$
Enlarging $C$ covers the finitely many small values of $M$.
Since $N\geq2$, the term $d=1$ occurs in $\Lcal_{<N}(M)$ and gives the lower
bound.
\end{proof}

\begin{proof}[Proof of Corollary~\ref{cor:uniform-q6}]
By definition, $\Lcal_{<N}(M)\leq\Lcal(M)$.
Substitute \eqref{eq:L-loglog} into Theorem~\ref{thm:main-q6} to obtain
\eqref{eq:loglog-uniform}.  Using \eqref{eq:L-prime-power} instead gives
\eqref{eq:prime-power-branch}.  Since $M\leq N^2$, we have
$1+\log\log(3M)\leq C\log(2N)$.  Applying this bound to both terms proves
\eqref{eq:common-log}.
\end{proof}

\begin{remark}\label{rem:two-slow-factors}
The logarithm in $N^2\log(2N)$ comes from primitive integer directions.
The exact Vinogradov subcount shows that its order is sharp at $M=N^2$.
The factor in $(N^4/M)\Lcal_{<N}(M)$ comes from modular directions with
$\gcd(2u,2v,M)>1$.  The rough estimate \eqref{eq:common-log} places a
logarithm on both terms and thus loses this distinction.
\end{remark}

\section{Lower bounds and sharp regimes}\label{sec:lower-bounds}

Recall that $\Ksix(N,M)$ is the best operator constant.
For $a=\ind_{I_N}$, the normalized sixth moment equals a solution count.
The passage from this count to the operator constant requires division by
the coefficient norm.

\begin{proposition}\label{prop:universal-lower}
For every $N\geq2$ and $N\leq M\leq N^2$,
$$
\Ksix(N,M)
\geq
\max\left\{\frac{J_{3,2}(N)}{N},\frac{N^4}{M}\right\}
\geq
c\left(N^2\log(2N)+\frac{N^4}{M}\right),
$$
where $c>0$ is absolute.
\end{proposition}

\begin{proof}
Take $a=\ind_{I_N}$.  By character orthogonality, the sixth moment counts
the sextuples in \eqref{eq:G-def}.  Every solution of the exact integer
equations satisfies these congruences for every $M$.  The translation
$x\mapsto x-1$ from $\{1,\ldots,N\}$ to $I_N$ preserves both equations.
For the quadratic equation, this follows from
$$
\sum_i(x_i-1)^2-\sum_i(y_i-1)^2
=\left(\sum_i x_i^2-\sum_i y_i^2\right)
-2\left(\sum_i x_i-\sum_i y_i\right).
$$
Thus
$$
\E|\cE_{N,M}\ind|^6\geq J_{3,2}(N).
$$
The single point $(A,B)=(0,0)$ also gives
$$
\E|\cE_{N,M}\ind|^6\geq\frac{N^6}{NM}=\frac{N^5}{M}.
$$
Divide by $\|\ind\|_6^6=N$ and use Lemma \ref{J}.
The inequality $\max\{X,Y\}\geq(X+Y)/2$ then gives the stated sum.
Decreasing the absolute constant $c$ covers the finitely many small values
of $N$.
\end{proof}

\begin{proposition}
\label{prop:crt}
If $N$ is odd and squarefree, then
\begin{equation}\label{eq:crt-exact}
\E_{A,B\bmod N}|\cE_{N,N}\ind(A,B)|^6
=N^4P(N),
\qquad
P(N):=\prod_{p\mid N}\left(1+\frac1p-\frac1{p^2}\right).
\end{equation}
In this case $P(N)=\Lcal(N)$.  If
$$
N_y:=\prod_{\substack{3\leq p\leq y \\ \text{$p$ prime}}}p,
$$
then
\begin{equation}\label{eq:primorial-growth}
P(N_y)\sim\log\log N_y.
\end{equation}
\end{proposition}

\begin{proof}
First let $p$ be an odd prime and define
$$
S_p(A,B):=\sum_{n\bmod p}e_p(An+Bn^2).
$$
When $B=0$, character orthogonality gives
$$
S_p(A,0)=
\begin{cases}
p,&A=0,\\
0,&A\neq0.
\end{cases}
$$
When $B\neq0$, completing the square and evaluating the quadratic Gauss sum
give
$$
|S_p(A,B)|=p^{1/2}
\qquad(A\bmod p).
$$
Hence the normalized local sixth moment is
\begin{equation}\label{eq:local-gauss-sixth}
\begin{aligned}
\E_{A,B\bmod p}|S_p(A,B)|^6
&=\frac1{p^2}\left\{p^6+p(p-1)p^3\right\}\\
&=p^4\left(1+\frac1p-\frac1{p^2}\right).
\end{aligned}
\end{equation}

For squarefree $N$, the Chinese remainder theorem decomposes the residue
system modulo $N$ into residue systems modulo $p\mid N$.  The unit factors
in the additive character induce bijections of the local pairs $(A,B)$.
Thus the normalized moments multiply.  Taking the product of
\eqref{eq:local-gauss-sixth} over $p\mid N$ proves \eqref{eq:crt-exact}, and
\eqref{eq:euler-product} gives $P(N)=\Lcal(N)$.

Finally,
$$
1+\frac1p-\frac1{p^2}
=\left(1-\frac1p\right)^{-1}
\left(1-\frac{2}{p^2}+\frac1{p^3}\right).
$$
Thus $P(N_y)$ is a Mertens product times an Euler product converging to a
positive constant, and so $P(N_y)\sim\log y$. Using Lemma \ref{estimate} gives $\log N_y\sim y$, proving
\eqref{eq:primorial-growth}.
\end{proof}

\begin{proposition}\label{prop:subgrid-lower}
Suppose that $M=LN$, where $1\leq L\leq N$ is an integer, and that $N$ is
odd and squarefree.  Then
$$
\E_{A\bmod N}\E_{B\bmod M}|\cE_{N,M}\ind(A,B)|^6
\geq\frac{N^5}{M}P(N),
$$
and consequently
\begin{equation}\label{eq:subgrid-K}
\Ksix(N,M)\geq\frac{N^4}{M}P(N).
\end{equation}
\end{proposition}

\begin{proof}
The $N$ points $B=Lb$, with $b\bmod N$, are distinct modulo $M$, and
$$
\cE_{N,M}\ind(A,Lb)
=\sum_{n=0}^{N-1}e_N(An)e_{LN}(Lbn^2)
=\cE_{N,N}\ind(A,b).
$$
It follows that
$$
\begin{aligned}
\E_{A\bmod N}\E_{B\bmod M}|\cE_{N,M}\ind|^6
&\geq\frac1{NM}\sum_{A,b\bmod N}|\cE_{N,N}\ind(A,b)|^6\\
&=\frac NM\E_{A,b\bmod N}|\cE_{N,N}\ind(A,b)|^6.
\end{aligned}
$$
Apply Proposition~\ref{prop:crt}, and then divide by
$\|\ind\|_6^6=N$ for \eqref{eq:subgrid-K}.
\end{proof}

\begin{corollary}
\label{cor:no-arithmetic-factor}
There is no absolute constant $C$ for which
\begin{equation}\label{eq:false-natural-bound}
\E|\cE_{N,M}a|^6
\leq
C\left\{N^2\log(2N)+\frac{N^4}{M}\right\}\|a\|_6^6
\end{equation}
holds for every $N\leq M\leq N^2$ and every complex sequence $a$.
\end{corollary}

\begin{proof}
Take $M=N=N_y$ and $a=\ind$.  Proposition~\ref{prop:crt} gives
$$
\E|\cE_{N,N}\ind|^6
=N^4P(N)\sim N^4\log\log N,
$$
whereas the right-hand side of \eqref{eq:false-natural-bound} is
$$
C\left(N^2\log(2N)+N^3\right)N=O(N^4).
$$
This is a contradiction as $N=N_y\to\infty$.
The obstruction also occurs at the transition scale.  Set
$$
L=\left\lfloor\frac{N}{\log N}\right\rfloor,
\qquad
M=LN\sim\frac{N^2}{\log N}.
$$
Proposition~\ref{prop:subgrid-lower} then gives a lower bound
$\gtrsim N^2\log N\log\log N$ for the operator constant, while
\eqref{eq:false-natural-bound} would give only $O(N^2\log N)$.
\end{proof}


\begin{theorem}
\label{thm:M-classification}
The following estimates hold.
\begin{enumerate}
\item If $M=p^\alpha$ is any prime power and $N<M\leq N^2$, then
\begin{equation}\label{eq:classification-prime-power-rectangle}
\Ksix(N,M)
\sim N^2\log(2N)+\frac{N^4}{M},
\end{equation}
and the single test sequence $a=\ind$ satisfies
\begin{equation}\label{eq:classification-prime-power-ones}
\E|\cE_{N,M}\ind|^6
\sim N^3\log(2N)+\frac{N^5}{M}.
\end{equation}

\item Let $N=N_y$ be the odd primorial in
\eqref{eq:primorial-growth}, and set
$$
L=\left\lfloor\frac{N}{\log N}\right\rfloor,
\qquad
M=LN\sim\frac{N^2}{\log N}.
$$
For sufficiently large $y$,
$$
\Ksix(N,M)\sim N^2\log N\log\log N,
$$
and
\begin{equation}\label{eq:classification-transition-primorial-ones}
\E|\cE_{N,M}\ind|^6
\sim N^3\log N\log\log N.
\end{equation}

\item For every $N$,
$$
\Ksix(N,N^2)\sim N^2\log(2N).
$$

\item There is an absolute $C_0$ such that, for all sufficiently large $N$,
\begin{equation}\label{eq:high-M-threshold}
M\geq
C_0\frac{N^2\bigl(1+\log\log(3N)\bigr)}{\log(2N)}
\end{equation}
implies
\begin{equation}\label{eq:classification-high}
\Ksix(N,M)\sim N^2\log(2N).
\end{equation}
The proportion of integers $M\in[N,N^2]$ satisfying
\eqref{eq:high-M-threshold} is
\begin{equation}\label{eq:high-M-density}
1-O\left(\frac{\log\log(3N)}{\log(2N)}\right).
\end{equation}
\end{enumerate}
\end{theorem}

\begin{proof}
For the first assertion, the upper bound in
\eqref{eq:classification-prime-power-rectangle} follows from
Corollary~\ref{cor:uniform-q6}.  Proposition~\ref{prop:universal-lower}
gives the lower bound.  Both tests in that proposition use $a=\ind$, so
$$
\E|\cE_{N,M}\ind|^6
\geq
\max\left\{J_{3,2}(N),\frac{N^5}{M}\right\}
\gtrsim N^3\log(2N)+\frac{N^5}{M}.
$$
Multiplying \eqref{eq:classification-prime-power-rectangle} by
$\|\ind\|_6^6=N$ gives the matching upper bound and proves
\eqref{eq:classification-prime-power-ones}.

For the second assertion, Proposition~\ref{prop:subgrid-lower} and
\eqref{eq:primorial-growth} give
$$
\Ksix(N,M)
\gtrsim\frac{N^4}{M}\log\log N
\sim N^2\log N\log\log N.
$$
The reverse bound follows from Theorem~\ref{thm:main-q6},
\eqref{eq:L-loglog}, and $M\sim N^2/\log N$.  The lower bound is realized
by $a=\ind$, and multiplying the upper bound for $\Ksix$ by $N$ proves
\eqref{eq:classification-transition-primorial-ones}.

For the third assertion, Proposition~\ref{prop:universal-lower} gives the
lower bound.  In Theorem~\ref{thm:main-q6}, the first term is $N^2\log(2N)$.
By \eqref{eq:L-loglog}, the second is $O(N^2\log\log(3N))$ and is absorbed
by the first.

Finally, \eqref{eq:high-M-threshold} and \eqref{eq:L-loglog} imply
$$
\frac{N^4}{M}\Lcal_{<N}(M)\leq CN^2\log(2N).
$$
Theorem~\ref{thm:main-q6} and Proposition~\ref{prop:universal-lower} prove
\eqref{eq:classification-high}.  There are $N^2-N+1$ admissible integers
$M$, and at most
$$
O\left(\frac{N^2\log\log(3N)}{\log(2N)}\right)
$$
fail \eqref{eq:high-M-threshold}.  This proves \eqref{eq:high-M-density}.
\end{proof}

\section{Finite aliasing and mixed-norm interpolation}
\label{sec:interpolation}

\begin{lemma}
\label{lem:alias}
Let $F:\T^2\to[0,\infty)$ be a trigonometric polynomial with
$$
F(x,t)=
\sum_{\substack{|u|\leq D_x\\|v|\leq D_t}}
\widehat F(u,v)e(ux+vt).
$$
Then
\begin{equation}\label{eq:alias}
\frac1{NM}\sum_{A=0}^{N-1}\sum_{B=0}^{M-1}F\left(\frac AN,\frac BM\right)
\leq
\left(1+\frac{2D_x}{N}\right)
\left(1+\frac{2D_t}{M}\right)
\int_{\T^2}F.
\end{equation}
\end{lemma}

\begin{proof}
Finite Fourier expansion and character orthogonality give
$$
\frac1{NM}\sum_{A,B}F\left(\frac AN,\frac BM\right)
=\sum_{\substack{N\mid u,\ M\mid v\\|u|\leq D_x,\ |v|\leq D_t}}\widehat F(u,v).
$$
The assumption $F\geq0$ does not imply that its Fourier coefficients are
nonnegative.  Instead, it gives
$$
\left|\widehat F(u,v)\right|\leq\int_{\T^2}F=\widehat F(0,0).
$$
Apply the triangle inequality to the alias sum.
Counting the multiples of $N$ and $M$ in the two frequency intervals proves
\eqref{eq:alias}.
\end{proof}

\begin{proposition}\label{prop:q2}
For every $\varepsilon>0$,
$$
\E|\cE_{N,M}a|^6\leq C_6\log N\left(1+\frac{N^2}{M}\right)\|a\|_{\ell^2}^6.
$$
\end{proposition}

\begin{proof}
Define
$$
S_a(x,t):=\sum_{n=0}^{N-1}a_ne(nx+n^2t),
\qquad
F(x,t):=|S_a(x,t)|^6.
$$
Since $F=S_a^3\overline{S_a}^{\,3}$, its Fourier support satisfies
$$
D_x\leq3(N-1),
\qquad
D_t\leq3(N-1)^2.
$$
Lemma~\ref{lem:alias} gives
\begin{equation}\label{eq:grid-torus}
\E|\cE_{N,M}a|^6
\leq
C\left(1+\frac{N^2}{M}\right)
\int_{\T^2}|S_a(x,t)|^6\dd x\dd t.
\end{equation}
To pass from the indexing $1\leq m\leq N$ in \eqref{eq:gly-input} to
$0\leq n<N$, set $m=n+1$.  The identity
$$
nx+n^2t=m(x-2t)+m^2t-x+t
$$
contributes a factor of modulus one after exponentiation.
The torus map $(x,t)\mapsto(x-2t,t)$ preserves Haar measure.  Hence
\eqref{eq:gly-input} gives
$$
\|S_a\|_{L^6(\T^2)}\leq C_6(\log N)^{1/6}\|a\|_{\ell^2}.
$$
Taking sixth powers and using \eqref{eq:grid-torus} proves the proposition.
\end{proof}

\begin{proof}[Proof of Theorem~\ref{thm:mixed-q}]
Regard $\cE_{N,M}$ as a finite-dimensional linear operator
$$
\cE_{N,M}:\ell^q(I_N)\longrightarrow L^6(G_{N,M},\mu_{N,M}).
$$
At the level of operator norms, Proposition~\ref{prop:q2} gives
$$
\|\cE_{N,M}\|_{\ell^2\to L^6}
\leq C_\varepsilon L_N^{1/6}(1+X)^{1/6},
$$
whereas Theorem~\ref{thm:main-q6} gives
$$
\|\cE_{N,M}\|_{\ell^6\to L^6}
\leq CN^{1/3}D_{N,M}^{1/6}.
$$

For $2\leq q<6$,
$$
\frac1q=\frac{1-\theta}{2}+\frac{\theta}{6},
\qquad
\theta=\frac32-\frac3q.
$$
Apply Riesz--Thorin interpolation to these two bounds and take sixth powers.
We obtain
$$
\E|\cE_{N,M}a|^6\leq C_{q,\varepsilon}N^{2\theta}L_N^{1-\theta}(1+X)^{1-\theta}
D_{N,M}^{\theta}\|a\|_{\ell^q}^6.
$$
Since $2\theta=3-6/q$, proves \eqref{eq:mixed-refined}.

Proposition~\ref{prop:euler-factor} and $M\leq N^2$ imply
$$
D_{N,M}\leq CL_N(1+X).
$$
Substituting this bound into \eqref{eq:mixed-refined} proves \eqref{eq:mixed-clean}.

Finally, if $q\geq6$, finite-dimensional H\"older gives
$$
\|a\|_{\ell^6}\leq N^{1/6-1/q}\|a\|_{\ell^q}.
$$
Substitution into Theorem \ref{thm:main-q6} yields
$$
\E|\cE_{N,M}a|^6
\leq CN^2D_{N,M}N^{1-6/q}\|a\|_{\ell^q}^6,
$$
which is \eqref{eq:mixed-qge6}.
\end{proof}

\section{A continuous cell consequence}\label{sec:continuous}

We define the half-open set
$$
\widetilde A_{N,M}
:=
[0,1)\times
\bigcup_{B=0}^{M-1}
\left[\frac BM,\frac BM+\frac1{N^2}\right).
$$
Since $M\leq N^2$, the intervals in the second coordinate have disjoint
interiors and
$$
\left|\widetilde A_{N,M}\right|=\frac{M}{N^2}.
$$

\begin{corollary}\label{cor:cells}
Under the assumptions of Theorem~\ref{thm:main-q6},
$$
\fint_{\widetilde A_{N,M}}
\left|\sum_{n=0}^{N-1}a_ne(nx+n^2t)\right|^6\dd x\dd t
\leq
C\left\{N^2\log(2N)+\frac{N^4}{M}\Lcal_{<N}(M)\right\}
\|a\|_{\ell^6}^6.
$$
The three mixed-norm conclusions of Theorem~\ref{thm:mixed-q} hold with the
same normalized average on the left-hand side.
\end{corollary}

\begin{proof}
Partition $[0,1)$ into $N$ half-open intervals and write, on each small
rectangle,
$$
x=\frac AN+v,\qquad t=\frac BM+u,\qquad 0\leq v<\frac1N,
\qquad
0\leq u<\frac1{N^2}.
$$
For fixed $(u,v)$, let
$$
a_n^{u,v}:=a_ne(nv+n^2u).
$$
Then $|a_n^{u,v}|=|a_n|$ and
$$
\sum_{n=0}^{N-1}a_ne\!\left(n\left(\frac AN+v\right)+n^2\left(\frac BM+u\right)\right)
=\cE_{N,M}a^{u,v}(A,B).
$$
The finite disjoint decomposition and Fubini's theorem give
$$
\begin{aligned}
&\fint_{\widetilde A_{N,M}}
\left|\sum_na_ne(nx+n^2t)\right|^6\dd x\dd t\\
&\qquad=
\frac{N^2}{M}
\int_0^{1/N}\int_0^{1/N^2}
\sum_{A=0}^{N-1}\sum_{B=0}^{M-1}
|\cE_{N,M}a^{u,v}(A,B)|^6\dd u\dd v.
\end{aligned}
$$
The conversion from the probability average contributes $NM$, the
$(u,v)$-rectangle has area $N^{-3}$, and
$$
\frac{N^2}{M}\cdot NM\cdot N^{-3}=1.
$$
Apply Theorem~\ref{thm:main-q6} for each fixed $(u,v)$.
Phase modulation preserves every $\ell^q$ norm.
The same argument with Theorem~\ref{thm:mixed-q} proves the mixed-norm bounds.
\end{proof}

\section{Open question}\label{sec:conclusion}

The remaining question in the strict rectangular range concerns the local
arithmetic factor.

\begin{problem}
For $N<M<N^2$. Theorem
\ref{thm:main-q6} gives
$$
\Ksix(N,M)\lesssim N^2\log(2N)+\frac{N^4}{M}\Lcal_{<N}(M),
$$
whereas Proposition~\ref{prop:universal-lower} gives
$$
\Ksix(N,M)
\gtrsim
N^2\log(2N)+\frac{N^4}{M}.
$$
The odd squarefree diagonal shows that the Euler factor can be sharp, and
Proposition \ref{prop:subgrid-lower} carries the obstruction to some
intermediate rectangles.  Is there an explicit arithmetic function $A(N,M)$,
defined from local prime-power data and not from $\Ksix$ itself, such that
$$
\Ksix(N,M)
\sim
N^2\log(2N)+\frac{N^4}{M}A(N,M)
$$
with absolute constants for all $N,M$?  If so, does $A(N,M)$ depend only on
$M$, or must it also involve $\gcd(N,M)$, the squarefree kernel, and local
Gauss moments?
\end{problem}


\begin{thebibliography}{99}

\bibitem{CushmanDemeterWu2026}
A.~Cushman, C.~Demeter and S.~Wu, Near optimal three-fold additive energy
bound for points on convex curves, arXiv:2608.12316, 2026.
\vspace{-0.3cm}

\bibitem{DFGGL26}
Y.~Deng, C.~Fan, S.~Guo, Z.~Guo and Y.~Luo, Sharp Strichartz estimates 
on the one-dimensional torus, arXiv:2609.18670, 2026.

\vspace{-0.3cm}

\bibitem{HickmanWright2018}
J.~Hickman and J.~Wright, The Fourier restriction and Kakeya problems over
rings of integers modulo $N$, Discrete Anal. (2018), Paper No.~11, 54 pp.

\vspace{-0.3cm}

\bibitem{Hu2026}
X. Hu, Complexity-sensitive additive energy and off-diagonal Young
inequalities on bounded-degree algebraic varieties, arXiv:2608.18956, 2026. 

\vspace{-0.3cm}

\bibitem{Johnsrude2024}
B.~Johnsrude, High-low analysis and small cap decoupling over
non-Archimedean local fields, arXiv:2409.09163, 2024.

\vspace{-0.3cm}

\bibitem{Johnsrude2025Sparse}
B.~Johnsrude, Sparse mean value estimates, algebraic number solution
counting, and non-Archimedean Fourier analysis, arXiv:2503.20015, 2025.

\vspace{-0.3cm}

\bibitem{KingsburyNeuschotz2025}
N.~Kingsbury-Neuschotz, On a restriction problem of Hickman and Wright
for the parabola over $\mathbb Z/N\mathbb Z$ for squarefree $N$,
arXiv:2509.09885, 2025.

\vspace{-0.3cm}

\bibitem{MontgomeryVaughan2006}
H.~L.~Montgomery and R.~C.~Vaughan, Multiplicative Number Theory I:
Classical Theory, Cambridge Stud. Adv. Math., 97 Cambridge University Press, Cambridge, 2007. xviii+552 pp.

\vspace{-0.3cm}

\bibitem{Rogovskaya1986}
N.~N.~Rogovskaya, An asymptotic formula for the number of solutions of
a system of equations, Diophantine approximations, Part II (Russian), Moskov. Gos. Univ., Moscow, 1986, pp. 78--84.

\vspace{-0.3cm}

\bibitem{VaughanWooley1997}
R.~C.~Vaughan and T.~D.~Wooley, A special case of Vinogradov's mean value
theorem, Acta Arith. 79 (1997), 193--204.

\end{thebibliography}

\bigskip \bigskip

\noindent Ziming Wang (Corresponding author) and Ziyi He

\medskip

\noindent Key Laboratory of Mathematics and Information Networks 
(Ministry of Education of China), School of Mathematical Sciences, 
Beijing University of Posts and Telecommunications, Beijing 100876, 
The People's Republic of China

\smallskip

\noindent{\it E-mails:} \texttt{zimingwang@bupt.edu.cn} (Z. Wang)\\
\noindent\hphantom{{\it{E-mails:}}} \texttt{ziyihe@bupt.edu.cn} (Z. He)
\end{document}
